\chapter{Conclusion and Future Work}
\label{ch:conclusion}

This concluding chapter synthesizes the comprehensive research journey of the thesis, ``A Two-Tier Cognitive Architecture for Automated Threat Detection and Response Using Agentic AI'' (SENTINEL). The chapter encapsulates the scientifically proven results mapped back to the core research questions, critically self-assesses the architectural and practical limitations within the scope of a master's thesis, and outlines future research directions for securing modern digital infrastructures using local, context-grounded Agentic AI.

\section{Synthesis of Achieved Research Results}
The synthesis of the achieved research results confirms the successful resolution of the three core research questions proposed at the outset of this study.

The first question (RQ1) concerned stateless wire-speed filtration. The design and implementation of Welford's online statistical algorithm at the transport tier (Tier-1) successfully addresses the challenges of high-volume telemetry processing and alert fatigue. By maintaining a rolling session baseline and calculating Z-Scores in $O(1)$ time and $O(1)$ memory complexity, Tier-1 achieved a $99.8\%$ benign-log filtration rate on the CSE-CIC-IDS2018 dataset. This substantial reduction ensures that only $0.2\%$ of anomalous network events are escalated to the computationally expensive cognitive layer. Furthermore, because Tier-1 maps active session baselines on a per-IP basis, it preserves the temporal event chains of low-and-slow Advanced Persistent Threats (APTs) without dropping the sequential context necessary for long-term correlation.

The second question (RQ2) addressed secure LLM triage and latency mitigation. The integration of a local, quantized Large Language Model (\textit{Gemma-2-9B-IT} Q6\_K) through a LangGraph-driven stateful FSM solves the dual challenges of operational latency and cognitive vulnerability. While the direct ingestion of raw logs by an LLM is bottlenecked by the quadratic complexity of self-attention, the SENTINEL architecture mitigates this by routing only pre-filtered anomalies and by utilizing an in-memory Semantic Cache (an LRU structure keyed on SHA-256 digests of event templates). This cache resolves recurring alerts in microseconds, whereas uncached anomalies average an inference latency of $4$--$6$ seconds per batch on consumer-grade hardware. On the security front, the Delimited Data Encapsulation layer (incorporating cryptographically generated dynamic hexadecimal nonces) together with the Output Sanitization filters proved highly effective at blocking log-substrate prompt injections (S1--S4) and delimiter smuggling, preventing malicious payloads from hijacking the agent's reasoning context.

The third question (RQ3) examined context grounding, memory, and statistical validation. The combination of the Dual-RAG retrieval system and the persistent Threat Memory SQLite database effectively grounds the agent's decisions and minimizes factual hallucination. Utilizing FAISS dense vector search and BM25 sparse keyword search merged via Reciprocal Rank Fusion (RRF, $k=60$), the RAG engine retrieves the relevant MITRE ATT\&CK techniques and NIST playbooks. This grounding was validated through a cross-family LLM-as-a-Judge evaluation, in which RAGAS-style metrics scored Faithfulness at $4.00$, Answer Relevancy at $4.63 \pm 0.80$, Context Recall at $3.99 \pm 0.07$, and Context Precision at $3.00 \pm 0.10$ on a $1$--$5$ scale. Long-term APT correlation is achieved through the Threat Memory schemas (\texttt{ip\_reputation}, \texttt{known\_entities}, \texttt{threat\_events}, and \texttt{apt\_indicators}), which escalate alert levels when a source IP exhibits suspicious activity spanning multiple days or kill-chain phases. Non-parametric statistical tests mathematically reinforce these findings: McNemar's test yielded a p-value of $1.0$ (indicating that Tier-2 adds explanatory power without altering the base F1 classification), and the Mann-Whitney U test yielded $U=27053.5$ with $p = 2.84 \times 10^{-17}$, confirming that the speedup achieved by the two-tier layout is statistically significant.

\section{Scientific and Practical Contributions}
This thesis delivers both scientific and practical contributions to the field of cybersecurity.

\subsection{Scientific Contributions}
The primary theoretical and scientific advancements of this work are threefold. The first is the formalization of a dual-tier cognitive security architecture that decouples transport-layer high-velocity data filtration (Tier-1) from cognitive-layer stateful reasoning (Tier-2), establishing a blueprint for combining deterministic mathematics with non-deterministic language models. The second is the development of a cryptographic isolation model for prompt templates---Delimited Data Encapsulation---that resolves the vulnerability of LLMs to untrusted log substrates without relying on fragile, compute-intensive classifier guardrails. The third is the introduction of a non-parametric evaluation framework for local cybersecurity agents, which employs cross-family LLM judges and RAGAS metrics to measure cognitive alignment, reasoning quality, and contextual grounding.

\subsection{Practical Contributions}
The empirical and engineering contributions of this work are likewise threefold. The first is a fully open-source, containerized prototype (SENTINEL), written in Python and orchestrated via Docker Compose, that integrates \texttt{llama.cpp} and \texttt{LangGraph} for rapid enterprise testing. The second is the validation of local air-gapped threat triage on consumer-grade hardware---an Intel Core i7-14700KF CPU, 32\,GB of DDR5 RAM, and a single NVIDIA RTX 4060 Ti GPU---ensuring data sovereignty and compliance with Zero-Trust guidelines for resource-constrained organizations. The third is an administrative dashboard (Streamlit UI) that bridges automation with human oversight through an interactive Human-in-the-Loop (HITL) queue and secures log files against database tampering by means of blockchain-like HMAC chaining.

\section{Existing Limitations}
Despite the system's demonstrated capabilities, several technical limitations define its operational boundaries. In terms of telemetry and format constraints, the implementation was evaluated primarily on network-flow telemetry and HTTP payload metadata; it lacks ready-to-use parser configurations for diverse unstructured enterprise logs---such as Windows Event Logs, Linux Syslogs, or CloudTrail---which would require additional Drain log-template extraction configurations. With respect to anomaly-queue saturation, although the Welford filter drops $99.8\%$ of benign events, a massive network attack such as high-rate scanning or denial of service, by causing a sudden spike in anomalies, would saturate the escalation queue and lead to VRAM constraints and queuing latency unless GPU compute is scaled. A further mitigation trade-off arises because, to prevent an accidental self-denial of service in which the LLM might block critical internal assets (such as Active Directory or database servers), the architecture relies on a human-in-the-loop mechanism to execute firewall-blocking actions; while this ensures safety, it limits the speed of autonomous real-time mitigation. Finally, with regard to online semantic poisoning, although vector-database checksums prevent offline file modification, the system remains vulnerable if an adversary can write to the knowledge directories through an authorized channel and thereby inject misleading playbook documents.

\section{Future Development Directions}
Several future development directions emerge from these limitations to guide subsequent research. A first direction is the adoption of multi-agent systems (MAS), decomposing the monolithic LangGraph agent into a cooperative network of specialized sub-agents---such as an IOC-extraction agent, a playbook-correlation agent, a firewall-rule-compiler agent, and a forensic-reporter agent---in order to enhance reasoning quality. A second direction is privacy-preserving federated learning, extending the architecture to a decentralized model in which multiple distributed SENTINEL nodes share locally trained model weights or reputation signals without exchanging raw, privacy-violating network-log payloads. A third direction is active reinforcement learning, integrating a local reinforcement-learning loop in which the agent updates its firewall-rule-generation parameters according to the explicit approvals or rejections issued by SOC analysts in the Streamlit HITL queue. A fourth direction is adversarial self-hardening, deploying a dedicated adversary-simulation agent inside the container network to probe SENTINEL automatically with prompt injections and evasion attacks, allowing the system to test and reinforce its own guardrail tokens dynamically.

\vspace{1cm}
\begin{flushright}
    \textit{Closing a modest endeavor, yet opening profound aspirations in the pursuit of securing the digital realm through the power of Agentic AI.}
\end{flushright}
